\section{Introduction}
%\subsection{Why QCS in CFD - QCFD}

Computational Fluid Dynamics (CFD) has become an indispensable third pillar in modern engineering sciences complementing theoretical and experimental analysis. Its broad applicability has led practitioners to constantly pushing the limits of numerical simulations for at least four decades. However, stagnation of Moore's law requires a radical rethinking of the way CFD codes and their underlying mathematical algorithms need to be designed in the future.

An emerging compute technology that has the potential to become a game changer for future CFD applications is quantum computing as it offers breakthrough solutions for the two major challenges of CFD today: large memory consumption and excessive need for computing power. The reader who is not familiar with the fundamental concepts of quantum computing is referred to the introductory lecture notes in this series for an introduction into quantum bits and gates, entanglement of quantum states, the superposition principle, quantum parallelism and the way of extracting information from a quantum computer by means of measurement. A more extensive description of these concepts and beyond can be found in \cite{mikeandike,DeWolf}. 

\normalem
In a nutshell, the advantage of quantum computers comes from their capability of encoding an exponentially large amount of data in linearly many quantum bits (qubits) and performing operations on all data simultaneously. This type of quantum parallelism is impossible with classical computers that either need to process data one by one, which results in an exponential growth in sequential run time, or duplicate the hardware resources and process multiple data in parallel, leading to an up to exponential growth in hardware resources. However, the advantage of quantum computing does not come for free. Extracting \emph{all} data from the quantum register requires an exponential number of runs and measurements, thereby foiling any potential quantum advantage. The art of designing quantum algorithms with practical advantage consists of reducing the number of necessary read-outs, e.g., by performing some post-processing of field data into scalar quantities of interest on the quantum computer itself.

CFD falls into the category of applications that have a good match with the capabilities and limitations of quantum computers, i.e., large amount of data to be worked with, high computational intensity, and, at the same time, users' primary interest in scalar-valued quantities of interest such as drag and lift coefficients, rather than the visual inspection of entire flow fields.

%\subsection{Literature review - general}
%\sout{Quantum computational fluid dynamics is a relatively new, but up and coming field with a plethora of applications and potential benefiting industries. Since computational fluid dynamics is constantly pushing the limits of modern super computers and computational power, it is a natural candidate to benefit from quantum computers.} 
This has led early adopters of this emerging compute technology to explore the possibilities of speeding up CFD applications through the use of quantum computers. In what follows, we give a brief overview over the state-of-the-art in quantum CFD literature.

%\sout{The idea is to exploit quantum parallelism to enable a speed-up for computational problems compared to classic fluid dynamics solutions. 
%Currently, there are quantum solutions for several types of computational fluid methods.}

In 2020 Gaitan published a quantum algorithm that can be used to solve the Navier-Stokes equations \cite{Gaitan2020}. In this article the author shows that there is a quantum speed-up for non-smooth flows and identifies a regime where the speed-up is quadratic over classical random algorithms.
Later, in 2021, Oz et al. adopted the algorithm presented in \cite{Gaitan2020} for solving partial differential equations to Burgers' equation \cite{Oz2021}. In this work the authors present solutions for flow problems with and without shock waves thereby achieving similar speed-up as for the Navier-Stokes equations. 
In a recent paper from 2022, Budinski suggests an alternative approach for solving the Navier-Stokes equations on a quantum computer by resorting to the stream-function vorticity formulation and adopting the lattice Boltzmann method \cite{Budinski2021}, thereby showing a proof of concept for yet another quantum method for solving the Navier-Stokes equations. 

A hybrid quantum classical reservoir computing model was suggested by Pfeffer et al. in 2022 \cite{Pfeffer2022}. With their method the authors are able to simulate nonlinear chaotic dynamics of Lorenz type models. They show that by using just a few highly entangled qubits they can achieve similar prediction and reconstruction capabilities as classical reservoirs using thousands of perceptrons. 

Another interesting approach that targets noisy intermediate quantum computers was proposed by Kyriienko et al. in \cite{Kyriienko2021}. In essence, a quantum neural network is trained to learn the solution values of a partial differential equation complemented by boundary and possibly initial values. The classical counterpart of this concept has become popular in recent years under the name physics-informed machine learning.

In 2021 Liu et al. presented a quantum algorithm for solving non-linear differential equations \cite{Liu2021}. The authors suggest to use Carleman linearization and subsequently perform time integration by the forward Euler method in combination with the quantum linear system algorithm by Harrow et al. \cite{Harrow2009} to find a solution. 

Another strategy for solving fluid flow problems with the aid of quantum computers is to resort to the Boltzmann equations and mimic the evolution of the particle distribution over time with a quantum circuit.
%\subsection{Quantum approaches for the Boltzmann equation}
In their seminal paper on the collisionless Boltzmann equation, Todorova and Steijl proposed quantum primitives for the streaming and specular reflection step, thereby marking the first quantum Boltzmann method (QBM) \cite{Todorova2020}.

In the same year, Budinski suggested a quantum lattice Boltzmann method for the D1Q2 and D2Q5 lattice structures that does include a simplified collision term, but cannot model the specular reflection step \cite{Budinski2020}. The collision term is realized using the linear approximation of unitary approach \cite{Low2019}, which requires a measurement at the end of each time step to check whether the previous computations have been meaningful or ended up in a so-called `orthogonal state of no interest'. In the former case, the algorithm can proceed to the next time step. In the latter case, the entire computation must be quashed and the simulation must be started from scratch. Post-selection of valid results after (partial) measurement as used in the approximation of unitary approach is a common practice in quantum algorithms but it becomes problematic if used repeatedly within a time-stepping loop. It is obvious that the probability of failure increases in a multiplicative manner with the number of timesteps. Therefore the amount of timesteps that can be feasibly modeled using this approach highly depends on the amplitudes of the `orthogonal states of no interest' and is not realistically implementable for an interesting amount of timesteps, unless the probability of failure can be shown to be very small. Unfortunately, the paper \cite{Budinski2020} does not provide any bound on this probability. 

In the subsequent paper \cite{Budinski2021}, Budinski proposed an implementation of the specular reflection step using the quantum linear approximation of unitary approach, which brings the same problems as before. Furthermore, the paper does not provide a clear procedure for decomposing the required unitary into known quantum gates which hinders its straightforward implementation. 

Another approach to solving the equations of computational fluid dynamics using a quantum computer is by making use of Carlemann linearization. In the papers \cite{Itani2022, Itani2023a,Sanavio2023} Succi et al. make use of the Carlemann linearization method to linearize the Boltzmann and other equations and subsequently solve this on a quantum computer. The downside of this method is that it leads to a system in which it is hard to run for more than one time step and which requires the decomposition of a large unitary.  

In this document we present our quantum Boltzmann methods. We do this by first presenting the collisionless quantum Boltzmann method we developed and explaining its building blocks. We subsequently present proofs that the commonly used encoding methods for quantum Boltzmann do not allow streaming and collision to both take place as a unitary operation. This implies that novel encoding methods are required to achieve a full quantum Boltzmann algorithm that does not require a 'stop-and-go' strategy. We subsequently present our so-called space-time encoding, which we developed to circumvent the cases where streaming and collision are not possible as one quantum unitary. Using our space-time encoding we are able to simulate multiple timesteps, whilst having the full operation be a quantum unitary and the operation is easily and inexpensively expressed as quantum gates without any decomposing of unitaries into quantum gates necessary. 

Our collisionless quantum Boltzmann method improves on the aforementioned quantum Boltzmann methods in various ways. Compared to \cite{Budinski2020, Budinski2021}, our approach enables a higher variety in the particle velocities. Moreover, our realization of the quantum specular reflection step is based on simple one- and two-qubit quantum gates, without the need to adopt the troublesome linear approximation of unitaries approach. This also allows us to ensure the straightforward implementability of our approach into quantum programming frameworks like Qiskit, without the need to decompose unknown unitaries or multi-control operations into native quantum gates. 



Our quantum space-time encoding marks the first quantum algorithm that can perform multiple timesteps of the lattice Boltzmann method as a full quantum unitary without the requirement of intermediate restarts. On top of that it sets itself apart from other quantum methods which allow for collision to take place in that it does not require the decomposition of exponentially large quantum unitary matrices. This makes our method much more realistic to be implemented on a quantum computer, as the decomposition of a quantum unitary can be exponentially expensive. 